% Bridge sections by /root/bridges.  All source locators are explicit in text.
% Bibliography keys and evidence URLs are listed in bridges_notes.md.

\section{Schubert $K$-theory, alternating positivity, and saturation}
\label{sec:bridge-k}

\paragraph{B3-KNOWN: the exact bridge.}
For the complex Grassmannian $\operatorname{Gr}(r,N)$, let
$b_\lambda=[\mathcal O_{\Omega_\lambda}]$ be its Schubert structure-sheaf
basis.  Write
\[
 b_\lambda b_\mu=\sum_\nu c^{K,\nu}_{\lambda\mu}b_\nu.
\]
Grothendieck polynomials are the representatives introduced by Lascoux and
Sch\"utzenberger; they are not polynomials written in the supplied
Grothendieck manuscript pages.  Buch's \emph{A Littlewood--Richardson rule
for the $K$-theory of Grassmannians}, Theorem~5.4 and Corollary~5.5,
express $(-1)^{|\nu|-|\lambda|-|\mu|}c^{K,\nu}_{\lambda\mu}$ as a
set-valued-tableau count.  His Introduction, immediately after the
definition of the product, identifies $c^{K,\nu}_{\lambda\mu}$ with the
ordinary Littlewood--Richardson coefficient when
$|\nu|=|\lambda|+|\mu|$, and says it is zero when
$|\nu|<|\lambda|+|\mu|$ \cite{buch2002}.
For a complex flag variety, Brion's \emph{Positivity in the Grothendieck
group of complex flag varieties}, Theorem~1, proves
\[
 (-1)^{\operatorname{codim}X_w-\operatorname{codim}Y}
 c_Y^w\geq0,
 \qquad [\mathcal O_Y]=\sum_w c_Y^w[\mathcal O_{X_w}],
\]
when $Y$ has rational singularities.  The product version follows from
Theorem~1 and Lemma~2, as explained directly after Theorem~1
\cite{brion2002}.  These hypotheses and the choice of basis matter.

The classical dilation equivalence says, for every integer $N\geq1$,
\[
 c_{N\lambda,N\mu}^{N\nu}>0
 \quad\Longleftrightarrow\quad c_{\lambda\mu}^{\nu}>0.
\]
An exact primary locator is Knutson--Tao, \emph{The honeycomb model of
$GL_n(\mathbb C)$ tensor products I}, \S6, Theorem~3 and its following
corollary converting honeycombs to tensor-product multiplicities,
with the forward semigroup property explained in its Introduction
\cite{kt1999}.  This is the additional integrality result in the
Littlewood--Richardson model, not a consequence of nonnegative
coefficients alone.

\subsection*{B3.1. The sign change and its precise computational scope}
\paragraph{Statement.}
Put $e_\lambda=(-1)^{|\lambda|}b_\lambda$.  The multiplication constants
in this basis are the nonnegative integers
$a_{\lambda\mu}^{\nu}=(-1)^{|\nu|-|\lambda|-|\mu|}
c^{K,\nu}_{\lambda\mu}$.  On the degree-equality locus
$|\nu|=|\lambda|+|\mu|$, an oracle for exact $K$-theoretic constants is
also an oracle for exact Littlewood--Richardson constants, with the same
three partitions as input.  In particular, changing the basis does not
produce a plethysm formula.
\paragraph{Proof.}
Substitute $b_\rho=(-1)^{|\rho|}e_\rho$ into the product:
\[
 e_\lambda e_\mu
 =\sum_\nu(-1)^{|\lambda|+|\mu|+|\nu|}
      c^{K,\nu}_{\lambda\mu}e_\nu
 =\sum_\nu a_{\lambda\mu}^{\nu}e_\nu.
\]
The two sign exponents differ by the even integer
$2(|\lambda|+|\mu|)$.  Nonnegativity is the cited Buch theorem.
On degree equality the cited identification gives
$a_{\lambda\mu}^{\nu}=c^{K,\nu}_{\lambda\mu}=c_{\lambda\mu}^{\nu}$,
so one oracle call proves the computational assertion.  Every operation
in this argument is a product or a sign change; it supplies no identity
for a Schur-functor composition.  This last sentence describes the scope
of the construction, not a formal independence theorem.\hfill$\square$

\subsection*{B3.2. Geometric alternating positivity does not preserve dilation support}
\paragraph{Statement.}
The sign-corrected $K$-theoretic coefficients do not satisfy the classical
dilation equivalence.  In fact
\[
 a_{(1),(1)}^{(2,1)}=1,
 \qquad a_{(2),(2)}^{(4,2)}=0.
\]
Both triples fit in the Schubert rectangle of $\operatorname{Gr}(2,6)$.
This refutes even the implication from nonvanishing at scale one to
nonvanishing at scale two, for these coefficients.  It does not refute
the separate, opposite implication on its own.
\paragraph{Proof.}
Use Lenart's Pieri formula in Buch, Corollary~4.7, together with
conjugation of partitions (Buch, Corollary~6.7; equivalently the
Grassmannian-duality identity in the Introduction).  In horizontal-strip
form it says
\[
 c^{K,\nu}_{\lambda,(p)}=
 (-1)^{|\nu/\lambda|-p}
 \binom{r(\nu/\lambda)-1}{|\nu/\lambda|-p}
\]
when $\nu/\lambda$ is a horizontal strip, where $r$ is its number of
nonempty rows, and the coefficient is zero otherwise.  For
$\lambda=(1)$, $p=1$, $\nu=(2,1)$, the two added boxes occupy two
different columns and two rows.  The coefficient is
$(-1)^1\binom11=-1$, also recorded in Buch, Example~5.6.
For $\lambda=(2)$, $p=2$, $\nu=(4,2)$, the four added boxes occupy
columns $3,4$ in the first row and columns $1,2$ in the second.  Hence
the strip is horizontal, $r=2$, and the coefficient is
$(-1)^2\binom12=0$.  Applying the defining sign correction proves both
values.\hfill$\square$

\paragraph{Obstruction gate for bridge 3.}
To transfer a $K$-theoretic computation into determinant-versus-permanent
GCT, one must identify a particular coordinate-ring multiplicity (or a
map detecting its failure to extend across the boundary) with the
specified $K$-theoretic coefficient, with all signs and degrees fixed.
Neither the name ``Grothendieck'' nor Brion's theorem provides that
identification.  B3.2 shows that even an available geometric positivity
theorem cannot simply be used as the Littlewood--Richardson saturation
step of the flip.  The supplied archive page cote~119, scan~217
(site page~216), discusses the $K$-invariant and $\lambda_i$ operations;
it contains neither the later Lascoux--Sch\"utzenberger polynomials nor
Brion's theorem.

\section{The metric inequality: what its positivity actually controls}
\label{sec:bridge-metric}

\paragraph{B4-KNOWN: dictionary and exact locators.}
For a real matrix $A=(a_{ij})$, define
\[
 B(A)=\max_{\varepsilon_i,\delta_j\in\{-1,1\}}
          \sum_{ij}a_{ij}\varepsilon_i\delta_j,
 \qquad
 S(A)=\max_{\|u_i\|=\|v_j\|=1}
          \sum_{ij}a_{ij}\langle u_i,v_j\rangle.
\]
The real Grothendieck inequality is $B(A)\leq S(A)\leq K_GB(A)$,
where one universal constant works for all dimensions.  The historical
primary locator is Grothendieck's 1953 \emph{R\'esum\'e de la th\'eorie
m\'etrique des produits tensoriels topologiques}, \S4, no.~2,
Theorems~1--2; its tensor-norm formulation is transported to the displayed
finite-matrix formulation in Alon--Naor, Introduction, programs~(1)--(2)
\cite{groth1953,alon-naor2006}.  Alon--Naor, Theorem~4.3, supplies
polynomial-time randomized rounding with expected objective at least
$2\log(1+\sqrt2)/\pi$ times the SDP objective.  Their Lemma~2.1 relates this
optimization to the cut norm.

Raghavendra--Steurer, \emph{Towards Computing the Grothendieck Constant},
Theorem~1.1, proves that a factor strictly better than $K_G$ for this
\emph{bipartite quadratic optimization problem} is NP-hard under the
Unique Games Conjecture.  This is not a threshold for every problem
involving tensors \cite{rs2009}.  In a two-party XOR game with a specified
target parity $f(i,j)$ and question distribution $\pi(i,j)$, set
$a_{ij}=\pi(i,j)(-1)^{f(i,j)}$.  Tsirelson's vector characterization
identifies the classical bias with $B(A)$ and quantum bias with $S(A)$.
For an accessible exact research-paper locator, see Cleve--H\o yer--Toner--Watrous,
\emph{Consequences and Limits of Nonlocal Strategies}, revised version
of 11 January 2010, Theorem~4 (Tsirelson), Proposition~5, and Theorem~7
\cite{chtw2010}.  Bias is twice winning probability minus one; this
qualification avoids conflating a bias ratio with a probability ratio.

\subsection*{B4.1. A zero-versus-positive norm realization already carries the decision problem}
\paragraph{Statement.}
For a rational real matrix $A$,
\[
 B(A)=0\quad\Longleftrightarrow\quad S(A)=0
 \quad\Longleftrightarrow\quad \|A\|_{\mathrm{cut}}=0
 \quad\Longleftrightarrow\quad A=0.
\]
Consequently, suppose a nonnegative integer function $m(z)$ has an
explicit polynomial-time construction of a polynomial-size rational
matrix $A_z$, with polynomial-bit entries, such that
$m(z)>0$ if and only if $B(A_z)>0$ (or the analogous condition with
$S$ or the cut norm).  Then deciding $m(z)>0$ is in deterministic
polynomial time, by inspecting entries alone.
\paragraph{Proof.}
The maximum of a bilinear form over two cubes is attained at vertices:
fix one cube variable at a time and move it to an endpoint that does not
decrease the objective.  Therefore $B(A)$ is also the maximum over
$x_i,y_j\in[-1,1]$.  For any entry $a_{ij}$, choose $x$ and $y$ supported
only at $i$ and $j$, with their signs chosen to make the objective
$|a_{ij}|$.  Thus $B(A)\geq\max_{ij}|a_{ij}|$.
Since scalar signs are admissible unit vectors, $S(A)\geq B(A)$.
A singleton rectangle gives
$\|A\|_{\mathrm{cut}}\geq\max_{ij}|a_{ij}|$.
All three quantities vanish for the zero matrix.  These observations
prove the equivalences.  The algorithm constructs $A_z$ and checks
whether any rational numerator is nonzero, taking polynomial time in
the asserted output size.\hfill$\square$

\subsection*{B4.2. A constant-factor estimate cannot generally decide a one-unit multiplicity gap}
\paragraph{Statement.}
Fix $\alpha>1$.  Knowing only a number $v$ satisfying
$c\leq v\leq\alpha c$ for an unknown positive integer $c$ does not
determine $c$, even when $c$ is promised to be one of two consecutive
integers.  Hence such an estimate by itself cannot certify every strict
comparison between two multiplicities.
\paragraph{Proof.}
Choose an integer $c\geq1/(\alpha-1)$.  The common output $v=c+1$
is permitted when the true value is $c$, since
$c+1\leq\alpha c$, and also when it is $c+1$.  An algorithm supplied
only with this guarantee and this output receives identical information
in the two cases.  Against a comparison value $c$, strict inequality
holds in one case and fails in the other.\hfill$\square$

\paragraph{Obstruction gate for bridge 4.}
The established passage is from a bilinear objective to a semidefinite
relaxation and a rounding guarantee.  GCT needs an explicit map from a
specified multiplicity or orbit-closure question to an optimization
instance with a preserved, adequately separated threshold.  B4.1 locates
the whole positivity difficulty in any proposed zero-preserving
matrix construction, and B4.2 records the precision needed for a
multiplicity obstruction.  None of the supplied cote/scan pages is the
1953 \emph{R\'esum\'e}.  No claim that RH or a subexponential SAT
algorithm resolves the Unique Games Conjecture is made here.

\section{Projective tensor norms and algebraic degenerations}
\label{sec:bridge-norm}

\paragraph{B5-KNOWN: exact objects.}
For finite-dimensional real or complex Hilbert spaces $V_1,\ldots,V_d$,
the projective tensor norm is
\[
 \|T\|_\pi=
 \inf_{T=\sum_j v_{1j}\otimes\cdots\otimes v_{dj}}
        \sum_j\prod_{i=1}^d\|v_{ij}\|.
\]
Normalizing each nonzero factor gives exactly the tensor nuclear norm
of Friedland--Lim, \emph{Nuclear Norm of Higher-Order Tensors},
equation~(1) and \S2.  Their Theorem~8.1 states NP-hardness for real
$d$-tensors for $d\geq3$; Theorem~8.10 supplies the complex statement
for $d\geq4$ in the arXiv v3 numbering.  These are norm-computation
statements, not tensor-rank lower bounds.  Base field matters for the
norm as well as for rank \cite{friedland-lim2018}.
Tensor rank counts decomposable summands without a norm bound.
Border rank asks membership in the Zariski closure of the tensors of
bounded rank.  For the matrix-multiplication tensor, a decomposition
$T_{\mathrm{MM}}=\sum_{j=1}^r u_j\otimes v_j\otimes w_j$ means precisely
$AB=\sum_{j=1}^r u_j(A)v_j(B)w_j$: substitution in the tensor identity
gives an algorithm using $r$ products of scalar linear forms.
The norm additionally depends on Hilbert metrics, which general linear
changes of bases do not preserve.

\subsection*{B5.1. Scalar changes rule out an unnormalized norm obstruction for a full linear orbit}
\paragraph{Statement.}
For $a\neq0$ and $T\neq0$,
\[
 R(aT)=R(T),\qquad \underline R(aT)=\underline R(T),
 \qquad \|aT\|_\pi=|a|\|T\|_\pi.
\]
Moreover, $aT$ belongs to the same
$\prod_i\operatorname{GL}(V_i)$-orbit as $T$.  In particular a
nontrivial absolute norm threshold cannot be constant on these orbits
or distinguish their scalar-invariant closure membership.
\paragraph{Proof.}
Multiplication of the first factor of every decomposable summand by
$a$ sends a rank-$r$ decomposition of $T$ to one of $aT$; multiplication
by $a^{-1}$ proves the reverse rank inequality.  The linear
automorphism $T\mapsto aT$ preserves each rank-at-most-$r$ locus, hence
its Zariski closure, proving the border-rank identity.  Applying the
same change to every decomposition in the definition of $\|\cdot\|_\pi$
gives $\|aT\|_\pi\leq|a|\|T\|_\pi$; the inverse change proves
equality.  Finally $(a\operatorname{id}_{V_1},\operatorname{id}_{V_2},
\ldots,\operatorname{id}_{V_d})$ is the required group element.
As $|a|$ ranges over positive reals, the norm crosses any fixed positive
threshold while the orbit is unchanged.\hfill$\square$

\subsection*{B5.2. The exact bounded-operator version that does carry an obstruction}
\paragraph{Statement.}
For linear maps $A_i:V_i\to W_i$ between finite-dimensional Hilbert
spaces,
\[
 \|(A_1\otimes\cdots\otimes A_d)T\|_\pi
 \leq\left(\prod_i\|A_i\|_{\mathrm{op}}\right)\|T\|_\pi.
\]
Therefore $\|S\|_\pi>\|T\|_\pi$ proves that $S$ is not in the
Euclidean closure of the tensors obtained from $T$ by maps satisfying
$\|A_i\|_{\mathrm{op}}\leq1$.  This statement does not impose a
restriction on a full $\operatorname{GL}$ orbit.
\paragraph{Proof.}
Apply the tensor product of the maps to an arbitrary decomposition of
$T$.  Each summand has projective cost at most
$\prod_i\|A_i\|_{\mathrm{op}}$ times its original cost, by the
defining operator-norm inequalities.  Sum and take the infimum.
When all maps are contractions, the resulting tensor has nuclear norm
at most $\|T\|_\pi$.  The norm is continuous in finite dimensions, so
the inequality persists under Euclidean limits.  B5.1 supplies the
unbounded scalar elements of the full general linear group, explaining
why that group is not covered.\hfill$\square$

\subsection*{B5.3. Even a natural normalization gives no uniform rank gap}
\paragraph{Statement.}
For every fixed order $d\geq2$ there are tensors $T_\epsilon$ of rank
and border rank two such that
\[
 \frac{\|T_\epsilon\|_\pi}{\|T_\epsilon\|_F}\longrightarrow1
 \quad(\epsilon\downarrow0),
\]
the same ratio as for a nonzero decomposable tensor.  All
$T_\epsilon$, $\epsilon>0$, lie in one general linear orbit.  Thus a
uniform positive additive gap in this ratio cannot separate rank one
from rank two.
\paragraph{Proof.}
Take
$T_\epsilon=(e_1\otimes e_1+\epsilon e_2\otimes e_2)\otimes
e_1\otimes\cdots\otimes e_1$, with every displayed $e_i$ unit and
orthogonal when indices differ.  It has a two-term decomposition.
Flattening between its first factor and all remaining factors has
matrix rank two for $\epsilon>0$.  Flattenings of a rank-one tensor
have matrix rank at most one, and their $2$-by-$2$ minors continue to
vanish on its Zariski closure.  This proves both asserted ranks.
The displayed decomposition gives
$\|T_\epsilon\|_\pi\leq1+\epsilon$.
The multilinear functional
\[
 L(v_1,\ldots,v_d)=
   \langle v_1,v_2\rangle_{\mathrm{bil}}
       \prod_{i=3}^d(v_i)_1
\]
has absolute value at most $\prod_i\|v_i\|$; over $\mathbb C$ the
first pairing is the bilinear coordinate pairing, for which the same
Cauchy--Schwarz bound holds.  Hence its linear extension satisfies
$|L(T)|\leq\|T\|_\pi$ by applying it to every decomposition.
Since $L(T_\epsilon)=1+\epsilon$, equality follows.  Orthogonality
gives $\|T_\epsilon\|_F=\sqrt{1+\epsilon^2}$, proving the limit.
The invertible diagonal map with entries $1,\epsilon$ on the first
factor sends $T_1$ to $T_\epsilon$.\hfill$\square$

\paragraph{Obstruction gate for bridge 5.}
A nuclear-norm bound can be used for a norm-constrained tensor model,
as B5.2 proves.  An algebraic orbit-closure obstruction requires a
quantity compatible with the allowed general linear action and with
degeneration, or an explicit theorem converting the bounded model into
the algebraic model with controlled loss.  Neither NP-hardness of
computing $\|\cdot\|_\pi$ nor a faster Boolean counting algorithm
supplies such a conversion.  The tensor operations in cote~2,
scan~17 are categorical pullback, pushforward, external tensor product,
and restriction along the diagonal; no Banach norm is specified there.

\section{Integral group schemes, Hilbert and Quot schemes, and fixed witnesses}
\label{sec:bridge-integral}

\paragraph{B7-KNOWN: what the foundations provide.}
SGA~3, Expos\'e~I (Demazure), \S4.2 describes affine group schemes by
their coordinate algebras, and \S4.5 constructs
\[
 \operatorname{GL}_{m,\mathbb Z}
 =\operatorname{Spec}\mathbb Z[t_{ij},\det(t)^{-1}].
\]
Thus integral tensor representations and their orbit maps can be
formed before choosing a characteristic \cite{sga3-I}.
Grothendieck, FGA, Bourbaki Expos\'e~221,
\emph{Les sch\'emas de Hilbert}, Theorems~3.1--3.2, represents the
functor of flat coherent quotients of a coherent sheaf on a projective
scheme over a noetherian base, with fixed Hilbert polynomial, by a
projective scheme.  The Hilbert scheme is the special case of a
quotient of the structure sheaf, immediately following Proposition~3.8.
Remark~3.9 explicitly reduces the universal construction to schemes
over $\mathbb Z$ and warns of components confined to a positive
characteristic \cite{fga221}.  These are the relevant exact historical
locators; they are not among the supplied cote/scan transcriptions.

\subsection*{B7.1. A characteristic-zero separation survives outside an explicit finite set}
\paragraph{Statement.}
Let $W$ be a finite free $\mathbb Z$-module with an algebraic action of
$G=\operatorname{GL}_{m,\mathbb Z}$, and let $x,y\in W$.
If $y$ does not belong to $\overline{G_{\overline{\mathbb Q}}x}$,
then there is $f\in\mathbb Z[W]$ with $f(y)\neq0$ whose restriction
to the orbit map $g\mapsto gx$ is identically zero over $\mathbb Z$.
For every prime $p\nmid f(y)$,
\[
 \overline y\notin
 \overline{G_{\overline{\mathbb F}_p}\overline x}.
\]
Consequently noncontainment of the corresponding orbit closures
survives at all these primes.  This assertion concerns one fixed pair
and one fixed witness; it asserts no uniform bound in a growing GCT
family.
\paragraph{Proof.}
Put $R=\mathbb Z[W]$ and
$A=\mathbb Z[t_{ij},\det(t)^{-1}]$.  The integral orbit map induces a
homomorphism $\phi_x:R\to A$.  Let
$I_{\mathbb Q}=\ker(\phi_x\otimes\mathbb Q)$.
Flatness of field extensions implies
$I_{\mathbb Q}\otimes\overline{\mathbb Q}
=\ker(\phi_x\otimes\overline{\mathbb Q})$.
This kernel is exactly the ideal of the orbit closure: a polynomial
vanishes on the orbit if and only if its pullback vanishes as a regular
function on the reduced group $G_{\overline{\mathbb Q}}$.
There is consequently $h$ in the extended ideal with $h(y)\neq0$.
Write $h$ as a finite $\overline{\mathbb Q}$-linear combination of
elements of $I_{\mathbb Q}$.  Evaluation is linear, so at least one of
those elements, say $f_0\in I_{\mathbb Q}$, has $f_0(y)\neq0$.
Clear its rational denominators to get $f\in R$.
Now $\phi_x(f)$ is zero in $A\otimes\mathbb Q$.
The ring $A$ is a localization of a polynomial ring over $\mathbb Z$,
hence torsion-free over $\mathbb Z$; therefore $\phi_x(f)=0$ in $A$.
Reduce this identity modulo $p$ and then extend to
$\overline{\mathbb F}_p$.  The reduced polynomial vanishes on the
reduced orbit, hence on its Zariski closure.  But it takes the nonzero
value $f(y)\bmod p$ at $\overline y$ if $p\nmid f(y)$, proving the
assertion.  Finally, for a group-stable closed set $X$,
$\overline{Gy}\subseteq X$ implies $y\in X$, and $y\in X$ implies
$\overline{Gy}\subseteq X$.  Apply this with the orbit closure of
$x$ in each field.\hfill$\square$

\subsection*{B7.2. Fixed-degree separation has an explicit rank criterion and stable reductions}
\paragraph{Statement.}
In B7.1, fix $d\geq0$.  There are an explicit integer matrix $M_d$
and an integer row $v_d$ such that, over any algebraically closed field
$k$, a polynomial of degree at most $d$ separates $y_k$ from
$\overline{G_kx_k}$ if and only if
\[
 \operatorname{rank}_k\begin{pmatrix}M_d\\v_d\end{pmatrix}
    >\operatorname{rank}_k M_d.                    \tag{B7.2}
\]
One can choose a nonzero integer $D_d$, as a product of at most two
nonzero rank-size minors, so that both ranks, and therefore this
separation condition, agree with characteristic zero whenever
$p\nmid D_d$.
\paragraph{Proof.}
Choose the monomials of degree at most $d$ as a basis of
$R_{\leq d}$.  Each of their pullbacks under $\phi_x$ is a regular
function on $G$, hence an integral polynomial in $t_{ij}$ divided by
a power of $\det(t)$.  There are finitely many basis monomials, so a
single power $\det(t)^K$ clears all these denominators.  Expand
$\det(t)^K\phi_x$ on this basis in ordinary monomials of $t_{ij}$,
using the finite union of all their supports as rows.  Its coefficient
matrix is $M_d$.  Let $v_d$ be evaluation of the original basis
monomials at $y$.  Over any field, the determinant polynomial is
nonzero and the polynomial ring is an integral domain.  Thus clearing
its power does not change whether a pullback is zero after localization.
The separating polynomials are exactly vectors $c$ such that
$M_dc=0$ and $v_dc\neq0$.

For a matrix over a field, a row vanishes on its kernel if and only if
it lies in its row space.  To see the nontrivial implication without an
extra hypothesis, a linear functional that vanishes on $\ker M_d$
factors through $\operatorname{im}M_d$; extend that functional to the
whole target by extending a basis.  This gives the desired row-space
expression.  Appending $v_d$ raises rank precisely when it fails to lie
in that row space.  This proves (B7.2).

For each of the two integer matrices choose a nonzero minor of size
its rational rank; if that rank is zero, use the empty minor $1$.
Let $D_d$ be their product.  A prime not dividing $D_d$ preserves these
two nonzero minors, so each reduced rank is at least its rational rank.
Every larger minor is the zero integer, so each reduced rank is at most
its rational rank.  Both ranks are equal as asserted.  Rank does not
change upon field extension, so the argument also applies to algebraic
closures.\hfill$\square$

\subsection*{B7.3. Bad-prime noncontainment need not lift}
\paragraph{Statement and proof.}
Let $\mathbb G_m$ act on $\mathbb A^1_{\mathbb Z}$ by scaling,
and take $x=2$, $y=1$.  Over $\overline{\mathbb Q}$ the orbit of $2$
is the nonzero locus, whose closure is all of $\mathbb A^1$; therefore
$y$ belongs to it.  In characteristic two, $x$ becomes zero, its orbit
closure is $\{0\}$, and $y=1$ lies outside.  This is a counterexample
to lifting arbitrary finite-characteristic noncontainment without a
good-reduction certificate.\hfill$\square$

\paragraph{Obstruction gate for bridge 7.}
B7.2 reduces a bounded-degree geometric separation to two concrete
linear algebra computations, and supplies a precise good-prime
certificate.  The unresolved GCT part is to obtain a sufficiently
small separating degree and a succinct, efficiently usable description
of the relevant matrices or representation summands for the growing
determinant/permanent family.  The monomial space has
$\binom{\operatorname{rank}W+d}{d}$ columns, so explicit expansion
alone does not supply a polynomial-time flip.  In positive
characteristic, one must also specify whether one is comparing formal
characters, Weyl filtrations, composition factors, or actual direct
summands; characteristic-zero semisimplicity cannot simply be reused.
Hilbert and Quot representability records families and boundaries, but
does not assert the needed degree bounds or multiplicity inequalities.
